% ECONOMICS_PROBLEM: {"id": "D81-16", "title": "Separate identification of beliefs and preferences", "jel_code": "D81", "source_id": "OP-0587"}
% FIRST_POSTED: 2026-10-11
% ATTEMPT_STATUS: solved
% ENTRY_KIND: research_attempt
% SUBMISSION: {"kind": "research_attempt", "category": "economics", "problem_id": "D81-16", "model": {"provider": "OpenAI", "version": "Codex, GPT-6 family", "version_uncertainty": "The system exposed the GPT-6 family; its exact serving revision was not exposed. No exact revision is inferred from the model folder.", "reasoning": "Not exposed by the session.", "generated_on": "2026-10-11"}, "other_models": "Separate Codex agents materially drafted and reviewed the proof. Their exact serving revisions and reasoning settings were not exposed; no other model identity is confirmed.", "tools": "File and Git/GitHub tools; primary-source web/PDF inspection; LaTeX compilation and page rendering. Python/Node site tests are publication checks. No Lean, computer algebra, or executed QE/LP implementation.", "target": "D81-16, Separate identification of beliefs and preferences: informative designs or model classes, sharp equivalence classes and identified policy functionals, extended to a declared ambiguity-sensitive criterion.", "scope": "Complete proof claim for the original informative-model-class target: finite consequences, invariant normalized utilities and beliefs, population probabilities and declared logit errors; calibrated max-min ambiguity with nonempty compact convex priors. All finite m,r>=1 and boundary/incompatibility cases are covered. No universal identification across unspecified error laws or ambiguity criteria is claimed.", "human_contributions": "JiHo Bae requested the problem range, readable and accurately cited proof, and submission. No significant human mathematical contribution or independent human proof checking is confirmed. The generating LLM is the primary mathematical author; JiHo Bae is credited as submitter.", "verification": "Separate Codex-agent proof, scope, boundary and primary-source reviews; symbolic example checks. These are model-assisted checks. Independent human verification: none confirmed. Formal verification: not performed. Document/site checks do not establish mathematical correctness.", "reproducibility": "Public source, bibliography, model-assisted review records and file hashes accompany the release. Original v1.0.0 preceded the new policy. Literal Prompt A was applied before the prospective v1.0.1 revision and Prompt B to its complete proposed source before submission; actual inputs and hashes were retained. No claim that these prompts preceded the original draft. Exact serving revision and full original generation transcript are unavailable."}
% FULL_SOLUTION_SCOPE: {"target": "D81-16, Separate identification of beliefs and preferences: informative designs or model classes, sharp equivalence classes and identified policy functionals, extended to a declared ambiguity-sensitive criterion.", "result": "Theorem 1: normalized matrix bijection, sharp observational fiber, structured-secant recognition and exact algebraic collision procedure. Theorem 2: constructive and minimal information dimension, including unknown scale. Theorems 3-4 and Proposition 1: entire compatible convex-prior family, sharp scalar and joint policy values, and whole-prior-set vertex criterion.", "coverage": "The original explicitly asks for informative designs or model classes. The manuscript specifies consequences, error law, reporting assumptions, utility restrictions and exclusions; gives injective designs and every noninjective fiber; identifies policy functionals; and separates calibrated preferences from convex and raw prior representations. All finite m,r>=1, zero prior coordinates, strict utilities, empty fibers/faces, J=0, L=0, known/unknown scale and lower-dimensional hulls are addressed.", "dependencies": "Basu-Pollack-Roy (2006), Sections 13.1 and 14.2: real-algebraic sampling and quantifier elimination over exact encoded coefficients. Rudin (1976), Theorem 9.24: C1 local inverse with nonsingular derivative for the dimension lower bound. Ben-Tal-Nemirovski, Optimization III (Spring 2023), Section 8.4.2, pp.235-239: polynomial bit complexity of rational linear programming on the bounded nonempty witness polytopes. Manski (2018), Sections II.A and VII.C: motivation, not the finite-experiment theorems.", "human_verification": "None confirmed. All mathematical reviews described here were separate Codex-agent checks, not independent human proof checking.", "unverified": "No independent human verification or Lean/formalization. The exact QE/LP procedures were characterized by proof and were not implemented or executed. Their cited hypotheses and symbolic boundary examples were checked by Codex agents. No missing mathematical obligation is asserted within the declared classes.", "novelty": "Checked the registered original statement, both earlier catalogue attempts, public repository/PR duplication and the cited primary sources. No exhaustive literature/equivalence search or established novelty claim is made."}
\documentclass[11pt]{article}
\usepackage[a4paper,margin=24mm]{geometry}
\usepackage[T1]{fontenc}
\usepackage{lmodern,microtype,amsmath,amssymb,amsthm,mathtools,xurl,hyperref}
\hypersetup{colorlinks=true,linkcolor=blue,urlcolor=blue,citecolor=blue}
\setlength{\parskip}{3pt}
\newtheorem{theorem}{Theorem}
\newtheorem{proposition}{Proposition}
\newcommand{\R}{\mathbb R}
\newcommand{\M}{\mathcal M}
\newcommand{\ind}{\mathbf1}
\newcommand{\conv}{\operatorname{conv}}
\title{Finite experiments separating beliefs, preferences and ambiguity}
\author{OpenAI Codex (GPT-6 family)\\\small Submitted by JiHo Bae}
\date{October 11, 2026}
\begin{document}
\maketitle

\paragraph{Submission provenance.}
The complete mathematical argument and its contribution disclosures are
retained from the \href{https://github.com/jbaelaw/belief-preference-identification-proof/blob/ce95354f5d52d298d0d3640f9af1de86bc36b266/paper/Proof.tex}{immutable public source} and
\href{https://github.com/jbaelaw/belief-preference-identification-proof/blob/ce95354f5d52d298d0d3640f9af1de86bc36b266/paper/Proof.pdf}{proof PDF}. This catalogue wrapper supplies identity and
submission metadata. The full-solution scope is the declared experimental
classes described in the manuscript. Community review remains pending.

\begin{abstract}
For a declared finite-consequence stochastic-choice model, we give an exact structural criterion for separate identification of beliefs and normalized utility. The sharp observational fiber is an affine slice of an explicitly defined rank-one matrix set. Exactly $m+r-2$ independent log odds suffice, and fewer cannot identify the unrestricted normalized family. A calibrated max--min extension gives the entire compatible convex-prior family, sharp policy-value intervals and a linear-programming description of joint policy values. A vertex criterion determines when the whole convex prior set is identified. The observations are population choice probabilities, rather than finitely many realized choices.
\end{abstract}

\section{A declared finite-choice model}
Catalogue problem D81-16 requests informative designs or model classes, sharp equivalence classes, identified policy functionals and an ambiguity-sensitive extension \cite{Catalogue}. We solve those requirements for the classes specified below. Manski discusses why beliefs and preferences can be confounded in choices and why elicited expectations need explicit interpretation \cite[\S\S II.A, VII.C]{Manski}. Our finite experiment results are not attributed to that survey.

There are $m\ge1$ states and consequences $c_0<\cdots<c_r$, with $r\ge1$. The utility family is
\[
 u_0=0<u_1<\cdots<u_r=1,
\]
where $u_j=u(c_j)$. Only utility on these consequences is a parameter. Beliefs satisfy $p\in\Delta_m$, including boundary priors. Fix ex ante a finite collection of supported environments. Each $z$ supplies a finite nonempty menu $A_z$ and a known consequence label $j_z(a,s)$. Put
\[
 V_{p,u}(a,z)=\sum_s p_su_{j_z(a,s)},\qquad
 \pi_z(a)=\frac{\exp(\lambda V_{p,u}(a,z))}
                   {\sum_{b\in A_z}\exp(\lambda V_{p,u}(b,z))},
 \quad \lambda>0.
\]
This is the declared choice-error law, with known common scale $\lambda$. The same fixed pair $(p,u)$ governs all observations; no mixture of unobserved decision-maker types is added. Menus change only the specified acts and consequences. Beliefs, utility and noise scale remain invariant, and environment assignment supplies no parameter-dependent information. Choices precede informative outcome disclosure. All comparisons below can instead be combined in one menu, avoiding learning between menus.

We observe population conditional choice probabilities. A finite experiment means finitely many supported menus; it does not mean finitely many sampled decisions recover continuous parameters exactly. No survey report is presumed truthful. If an exact reported vector $e$ is used, imposing $p=e$ requires an additional truthful-reporting assumption; otherwise its specified response law must be included in the observation model.

Choose a reference action $a_z^0$ in each menu and let
\[
 N=\sum_z(|A_z|-1),\qquad
 \ell_{z,a}=\lambda^{-1}\log\frac{\pi_z(a)}{\pi_z(a_z^0)}
 \quad(a\ne a_z^0).
\]
Model probabilities are strictly positive. Zero observed probabilities are incompatible with this finite-scale model.

\section{A structural criterion and sharp fibers}
Let $W\in\R^{m\times r}$, with columns indexed by $j=1,\ldots,r$. Define $\M$ by
\begin{align}
 W_{sj}&\ge0,\qquad \sum_sW_{sr}=1,\label{normalization}\\
 0&<\sum_sW_{s1}<\cdots<\sum_sW_{sr}=1,\label{order}\\
 W_{sj}&=W_{sr}\sum_tW_{tj}\qquad(j<r).\label{rankone}
\end{align}
For $r=1$, \eqref{order} says only $0<1$, and \eqref{rankone} is empty. Construct the integer matrix
\[
 D_{(z,a),(s,j)}=\ind\{j_z(a,s)=j\}
                    -\ind\{j_z(a_z^0,s)=j\},\qquad j\ge1.
\]

\begin{theorem}\label{identification}
The map $W_{sj}=p_su_j$ is a bijection from the normalized parameter family to $\M$. Its inverse is
\[
 p_s=W_{sr},\qquad u_j=\sum_sW_{sj}.
\]
The sharp equivalence class at observation $\ell$ is
\begin{equation}\label{fiber}
 \mathcal F(\ell)=\{W\in\M:D\operatorname{vec}W=\ell\},
\end{equation}
with parameters recovered by that inverse. The experiment is globally identifying if and only if
\begin{equation}\label{secant}
 \ker D\cap(\M-\M)=\{0\}.
\end{equation}
There is a terminating exact recognition procedure for this property from the finite consequence maps.
\end{theorem}
\begin{proof}
Generated matrices satisfy \eqref{normalization}--\eqref{rankone}. Conversely their last column is a probability vector, their column sums are a strictly ordered normalized utility, and \eqref{rankone} gives their product. This includes zero prior coordinates: the corresponding rows vanish, while column sums still recover utility.

The choice formula gives exactly $D\operatorname{vec}W=\ell$. Conversely this equality recovers every within-menu probability ratio; normalization recovers every probability. Thus every retained matrix realizes the observations, and no other parameter pair does. An empty fiber reports incompatibility.

Two parameter pairs coincide exactly when their matrices coincide. Distinct pairs give equal observations exactly when their difference is a nonzero structured secant in $\ker D$, proving \eqref{secant}. To recognize the property, test the finite polynomial system
\[
 W,\widetilde W\in\M,\qquad
 D\operatorname{vec}(W-\widetilde W)=0,\qquad
 \sum_{s,j}(W_{sj}-\widetilde W_{sj})^2>0.
\]
Real quantifier elimination decides feasibility \cite[Sections 13.1, 14.2]{BPR}. A nonempty system defined over the rationals has a real-algebraic sample point, giving an explicit collision. An empty system proves identification. The test uses the actual strict parameter set, rather than its closure or an unrestricted rank-two relaxation. No polynomial-time bound is claimed for this general recognition procedure.
\end{proof}

The criterion depends on primitive consequence variation through $D$ and on utility restrictions through $\M$. Any declared semialgebraic utility restrictions can be added to \eqref{order} throughout. A separate label $\theta$ is identified only if its map to the normalized finite utility vector is injective.

For a policy act on the same consequences, expected utility and its contrasts are known linear functions $L(W)$. Their exact sharp identified set is
\[
 \{L(W):W\in\mathcal F(\ell)\}.
\]
This is a semialgebraic image, computable by quantifier elimination when transformed data and target coefficients have an exact rational or real-algebraic representation. General semialgebraic target formulas can be treated in the same way. Global recognition above needs only rational $D$; effective computation at observed data additionally needs this coefficient convention. Arbitrary real probabilities do not automatically provide an effective encoding of their logarithms.

On a nonempty fiber, a linear target is identified exactly when its supremum and infimum agree. Extrema need not be attained because utility inequalities are strict, and the sharp set need not be an interval. For example, let $m=r=2$, write $p=p_1$, $u=u_1$, and observe acts $(c_1,c_0)$ and $(c_r,c_1)$ against constant $c_0$. Transformed values $3/16$ and $13/16$ give
\[
 pu=3/16,\quad p+(1-p)u=13/16,
 \qquad (p,u)\in\{(1/4,3/4),(3/4,1/4)\}.
\]
An endpoint state-one bet then has sharp value set $\{1/4,3/4\}$, while the observed first act has identified value $3/16$. The utility restriction $u>1/2$ selects the first pair. For a nonattained bound, observing $(c_1,c_r)$ against zero with value $1/2$ gives $p(1-u)=1/2$ and the sharp prior set $p\in(1/2,1]$.

\section{The smallest information dimension}
\begin{theorem}\label{design}
For the unrestricted normalized finite utility family and known $\lambda$, the smallest number of independent log odds of a globally identifying finite experiment is
\[
 d=m+r-2.
\]
If the positive common scale is also unknown, joint identification of $(\lambda,p,u)$ requires and admits $d+1$ such coordinates.
\end{theorem}
\begin{proof}
For each $s<m$, an act paying $c_r$ only in state $s$ has value $p_s$. For each $j<r$, a sure act paying $c_j$ has value $u_j$. Compare these $d$ acts with constant $c_0$, either in $d$ binary menus or one menu with $d+1$ actions. Log odds divided by $\lambda$ recover all these values, and $p_m=1-\sum_{s<m}p_s$. This works at boundary priors too.

For the lower bound we use the following elementary fact. A $C^1$ map $F$ from a nonempty open subset of $\R^d$ to $\R^N$, with $N<d$, cannot be injective. Let $\rho$ be its largest derivative rank. If $\rho=0$, the map is constant on a small ball. Otherwise reorder coordinates so a nonzero $\rho$-minor comes first. Near that point rank is constantly $\rho$, and
\[
 G(x)=(F_1(x),\ldots,F_\rho(x),x_{\rho+1},\ldots,x_d)
\]
is a local diffeomorphism by the inverse function theorem \cite[Theorem 9.24]{Rudin}: $G$ is $C^1$ on an open neighborhood and its derivative there is invertible at the chosen point. In these coordinates the first $\rho$ outputs are the first $\rho$ coordinates. Rank $\rho$ forces every other output derivative in the remaining coordinates to vanish. On a small product box the outputs are independent of those $d-\rho>0$ coordinates, giving a collision.

Strictly positive priors and strictly ordered utilities form a nonempty open set in the $d$ independent coordinates $(p_1,\ldots,p_{m-1},u_1,\ldots,u_{r-1})$. Their log-odds map $D\operatorname{vec}(pu^\top)$ is polynomial. Fewer than $d$ outputs cannot be injective by the fact just proved. If $d=0$, then $m=r=1$, the parameter pair is unique and no informative observation is needed.

For an unknown common scale, add sure $c_r$ against sure $c_0$. Their raw log odds are $\lambda$, after which the previous reconstruction applies. The lower bound applies to the $C^1$ raw-log-odds map $\lambda D\operatorname{vec}(pu^\top)$ on an open domain of dimension $d+1$. This proves joint minimality, including when $d=0$.
\end{proof}

Thus full column rank of $D$ is sufficient but unnecessary: the identifying construction uses $d<mr$ observations to recover the structured matrix. The minimum concerns information coordinates, not the number of arbitrary-size menus or sampled decisions.

\section{Calibrated ambiguity and its sharp prior family}
We now declare a separate max--min model. Replace the single prior by a nonempty compact convex set $K\subseteq\Delta_m$ and use
\[
 V_K(a)=\min_{p\in K}\sum_s p_su_{j(a,s)}.
\]
Choices obey the same logit law applied to these values. The same $K,u,\lambda$ remain fixed across menus under the earlier exclusions. Sure acts against $c_0$ identify every $u_j$ independently of $K$. The endpoint comparison also identifies an unknown common scale if necessary.

After this calibration, observe finitely many acts $a_j$ against constant $c_0$. They may share one anchored menu. Write their known state-utility vectors and observed values as
\[
 d_j=(u_{j(a_j,1)},\ldots,u_{j(a_j,m)}),\qquad
 v_j=\lambda^{-1}\log\frac{\pi(a_j)}{\pi(c_0)}.
\]
For $J\ge1$ such acts, define the compact polytope and its observed faces
\begin{equation}\label{faces}
 H=\{p\in\Delta_m:p\cdot d_j\ge v_j\ \forall j\},\qquad
 F_j=H\cap\{p:p\cdot d_j=v_j\}.
\end{equation}

\begin{theorem}\label{priorfamily}
The calibrated data are compatible if and only if $H$ and every $F_j$ are nonempty. Their entire sharp convex-prior family is
\begin{equation}\label{family}
 \mathcal K(v)=\{K:\varnothing\ne K\subseteq H,\ K\text{ compact convex},
                         \ K\cap F_j\ne\varnothing\ \forall j\}.
\end{equation}
\end{theorem}
\begin{proof}
Every compatible $K$ lies in $H$, and each observed minimum is attained in $K\cap F_j$. Conversely the conditions in \eqref{family} make every observed minimum exactly $v_j$, hence reproduce every anchored logit observation. If the faces are nonempty, their witness points' convex hull is compatible, as is $H$ itself. If a required face is empty, no compatible prior set exists. Utility-calibration failures are likewise reported as incompatibility before applying this theorem.
\end{proof}

\section{Policy values, joint contrasts and prior-set identification}
\begin{theorem}\label{policies}
For compatible data with $J\ge1$, the sharp set of lower policy utility $\min_{p\in K}p\cdot h$ is exactly
\begin{equation}\label{interval}
 \left[\min_{p\in H}p\cdot h,\quad
              \min_j\max_{p\in F_j}p\cdot h\right].
\end{equation}
For vectors $h_1,\ldots,h_L$, the entire joint identified set of their lower utilities $w_1,\ldots,w_L$ is the projection of this finite linear system: introduce $p^j\in H$ for each observed face and $q^\ell\in H$ for each policy, and require
\begin{align}
 p^j\cdot d_j&=v_j,& q^\ell\cdot h_\ell&=w_\ell,\label{witness-equalities}\\
 p^j\cdot h_k&\ge w_k,&q^\ell\cdot h_k&\ge w_k
                           &&\text{for every applicable index}.\label{witness-inequalities}
\end{align}
This joint set is a compact polytope. Every retained value vector has a witnessing prior set that is the convex hull of at most $J+L$ points. All linear policy contrasts have sharp extrema obtained from this same system; with rational calibrated coefficients and transformed data, this is rational linear programming.
\end{theorem}
\begin{proof}
Containment $K\subseteq H$ gives the lower endpoint of \eqref{interval}, attained at $K=H$. Since $K$ meets every $F_j$, its minimum under $h$ is at most each face maximum. Choose a maximizer $p_j$ of $h$ on each $F_j$. Their convex hull is compatible and attains the upper endpoint. For any intermediate $w$, choose $p_0\in H$ with $p_0\cdot h=w$, possible by convexity since the upper endpoint is at most $\max_Hp\cdot h$. Adding $p_0$ to those face maximizers gives a compatible convex hull whose minimum is exactly $w$. Thus every interval point is attainable.

For the joint result, any compatible $K$ supplies minimizers $p^j$ for the observed acts and $q^\ell$ for the policies. They satisfy \eqref{witness-equalities}--\eqref{witness-inequalities}. Conversely take the convex hull of all such witnesses. The inequalities bound each policy value below by $w_\ell$, its own witness attains equality, and every original observed face is still met. This gives exactly the requested joint values and proves sharpness.

The witness constraints are linear. Witnesses lie in a simplex, and the policy equalities bound every $w_\ell$ between the smallest and largest coordinates of $h_\ell$. The feasible system is therefore a compact polytope, as is its projection. Optimizing any linear contrast over it gives attained sharp endpoints. A scalar contrast is identified exactly when these endpoints agree. Marginal intervals alone cannot generally be subtracted to obtain a sharp contrast interval.
\end{proof}

All scalar extrema in \eqref{interval} are ordinary linear programs on compact polytopes. For rational calibrated coefficients and transformed data, these and the joint witness programs have polynomial bit complexity by rational linear programming \cite[Section 8.4.2]{BenTalNemirovski}. The coefficients are explicitly encoded rationals and the nonempty feasible polytopes are bounded, so the cited bit-complexity result applies. With arbitrary real coefficients the mathematical characterizations remain valid, while effective computation requires an exact input convention. No polynomial bound for whole-prior-set vertex enumeration is implied.

\begin{proposition}\label{vertices}
For compatible data with $J\ge1$, the entire convex prior set is identified if and only if every vertex of $H$ is the unique point of some observed face $F_j$. If it is identified, it equals $H$.
\end{proposition}
\begin{proof}
Singleton faces force their points into every compatible $K$. If they force every vertex, convexity gives $H\subseteq K\subseteq H$. Conversely, if a vertex $v$ is not uniquely forced by any face, choose $p_j\in F_j\setminus\{v\}$ for every $j$. Their convex hull is compatible but omits $v$, since an extreme point cannot be a convex combination of other points of $H$. It differs from the compatible set $H$, proving nonidentification. This includes lower-dimensional polytopes.
\end{proof}

When $J=0$, the sharp family is all nonempty compact convex subsets of $\Delta_m$. A scalar policy has sharp interval $[\min_{\Delta_m}p\cdot h,\max_{\Delta_m}p\cdot h]$, and the whole prior set is identified exactly when $m=1$. For $J=0,L\ge1$, the joint system uses only the policy witnesses. For $J=L=0$, any singleton prior is a witness and there are no policy coordinates to constrain.

If the primitive representation is a raw nonempty set $\Pi$, then
\[
 \inf_{p\in\Pi}p\cdot h
       =\min_{p\in\overline{\conv}\Pi}p\cdot h.
\]
Its sharp family consists of raw sets whose closed convex hull belongs to \eqref{family}. Different hulls can agree on finitely many observed directions. The vertex criterion identifies the convex hull, not an unrestricted raw representation. A nonsingleton hull and that hull with a nonextreme point removed are distinct raw sets with identical evaluations; a unique singleton hull does identify the raw set.

For a concrete example, take two states and endpoint consequences, and observe one state-one bet with value $a\in(0,1)$. Writing $p$ for the state-one probability gives $H=[a,1]$ and $F_1=\{a\}$. The sharp family is $K=[a,t]$, $t\in[a,1]$. The observed bet has identified lower utility $a$, but the state-two policy has sharp interval $[0,1-a]$, since its lower utility is $1-t$. The unforced vertex one explains why the whole prior set remains unidentified.

The theorems establish informative experimental designs, sharp equivalence classes and policy identification throughout the declared classes, including calibrated ambiguity. The error law and invariance conditions are maintained assumptions; utility outside the consequence set is not a parameter of these models.

\section{Scope, contributions and verification}
The original problem expressly asks for informative designs or model classes, rather than universal identification of unrestricted beliefs and preferences. Theorem~\ref{identification} gives a primitive consequence-matrix criterion, an exact recognition procedure and every sharp expected-utility fiber. Theorem~\ref{design} constructs an identifying design and proves its minimal information dimension. Theorems~\ref{priorfamily} and~\ref{policies}, with Proposition~\ref{vertices}, settle the declared ambiguity extension: all compatible convex prior sets, all joint policy values and whole-set identification. These statements cover every finite $m,r\ge1$, boundary priors, incompatible observations, nonattained bounds caused by strict utility restrictions, $J=0$ or $L=0$, and raw prior representations as specified above. No truthful survey response is assumed; the consequence and menu variations, invariances and choice-error law are explicit. This is a complete proof claim for the model-class target. It does not claim identification under every possible error law or ambiguity criterion.

The mathematics was generated and checked by OpenAI Codex agents on October 11, 2026. The exposed system identity is the GPT-6 family; the exact serving revision and reasoning configuration were not exposed. No other model identity can be confirmed. JiHo Bae requested the problem range, readability, source accuracy and publication, and is credited as submitter. No significant human mathematical contribution or independent human proof checking has been confirmed. Separate Codex agents reviewed the proofs, original-target coverage, boundary cases, citations and presentation; those are model-assisted reviews. The examples were checked symbolically. The exact quantifier-elimination and linear-programming procedures are proved characterizations; no implementation of those procedures was executed. No Lean formalization or computer algebra calculation is claimed. LaTeX compilation and PDF inspection check document production only.

The first public version preceded the repository's new contribution policy. This revision preserves its mathematical results, corrects the attribution and adds the required declarations and dependency citations. The original commit, tag and release assets remain available. The policy's literal prompt A was applied before preparing this revision, and prompt B was applied to its complete proposed source before submission; the audit records retain the actual inputs and source hashes. This chronology does not claim that those prompts preceded the original draft. The Manski source and the cited mathematical dependencies were checked, as were the catalogue's two earlier attempts. No exhaustive literature or novelty search, independent human verification or formal proof checking has been performed. The remaining verification task is independent expert review, not a missing step asserted within the mathematical argument.

\clearpage
\section{Completion Estimate}
\noindent\textbf{COMPLETION: 100\%}
The constructive designs, exact recognition and sharp fibers resolve the original informative-model-class target; the compatible-prior family, joint policy values and vertex criterion resolve its ambiguity extension. All cases within those explicitly declared classes are covered. This percentage states claimed completion of that target, not confidence or independent verification.

\begin{thebibliography}{9}
\raggedright
\bibitem{Catalogue} Problem Hunting with LLMs.
\emph{D81-16: Separate identification of beliefs and preferences}.
Economics catalogue, source ID OP-0587, statement at commit
\texttt{738965af403ff4375a58450fd971c43a817843bc}.
\url{https://github.com/mehmetmars7/Erdosproblems-llm-hunter/blob/738965af403ff4375a58450fd971c43a817843bc/attacks/open_problems/economics/statements/D81-16.tex}.
\bibitem{Manski} Charles F. Manski.
\emph{Survey Measurement of Probabilistic Macroeconomic Expectations: Progress and Promise}.
NBER Macroeconomics Annual 32 (2018), 411--471, Sections II.A and VII.C.
\url{https://doi.org/10.1086/696061}.
\bibitem{BPR} Saugata Basu, Richard Pollack and Marie-Fran\c{c}oise Roy.
\emph{Algorithms in Real Algebraic Geometry}, second edition.
Springer, 2006, Sections 13.1 and 14.2.
\url{https://doi.org/10.1007/3-540-33099-2}.
\bibitem{Rudin} Walter Rudin.
\emph{Principles of Mathematical Analysis}, third edition.
McGraw--Hill, 1976, Theorem 9.24.
\url{https://www.lehman.edu/faculty/rbettiol/lehman_teaching/2020mat320/baby_Rudin.pdf}.
\bibitem{BenTalNemirovski} Aharon Ben-Tal and Arkadi Nemirovski.
\emph{Lecture Notes: Optimization III}, ISYE 6663, Spring Semester 2023.
Section 8.4.2, ``Khachiyan's Theorem,'' pp.\ 235--239.
\url{https://www2.isye.gatech.edu/~nemirovs/OPTIIILN2023Spring.pdf}.
\end{thebibliography}
\end{document}
